\begin{tikzpicture}[
  x=1mm,y=1mm,
  font=\figurefont\footnotesize,
  panel/.style={draw=BookRule,fill=NNPanel,line width=.85pt,rounded corners=1.5mm},
  input/.style={draw=NNInput,fill=NNInput!16,line width=1pt,minimum width=7mm,minimum height=7mm,inner sep=1mm,align=center},
  compute/.style={draw=NNHidden,fill=NNHidden!16,line width=1pt,minimum width=10mm,minimum height=8mm,inner sep=1mm,align=center,rounded corners=1mm},
  output/.style={draw=NNOutput,fill=NNOutput!16,line width=1pt,minimum width=8mm,minimum height=8mm,inner sep=1mm,align=center},
  loss/.style={draw=NNGradient,fill=NNGradient!12,line width=1pt,minimum width=10mm,minimum height=8mm,inner sep=1mm,align=center,rounded corners=1mm},
  forward/.style={knowledge arrow,draw=NNWire,line width=1pt},
  gradient/.style={knowledge arrow,draw=NNGradient,line width=1.1pt,dashed,rounded corners=1mm},
  gradient stop/.style={-{Bar[width=2.4mm]},draw=NNGradient,line width=1.1pt,dashed,rounded corners=1mm},
  ema/.style={knowledge arrow,draw=BookTeal,line width=1pt,dash dot},
  note/.style={font=\figurefont\scriptsize,text=BookMuted,align=center}
]
  \draw[panel] (.5,103) rectangle (168.5,153);
  \draw[panel] (.5,48) rectangle (168.5,98);
  \draw[panel] (.5,-10) rectangle (168.5,42);
  \node[font=\figurefont\bfseries\footnotesize,text=BookInk,anchor=west] at (5,148) {(a) 对比判别：SimCLR};
  \node[note,anchor=east] at (163,148) {两路共享编码器与投影器};
  \node[input] (s-v1) at (10,132) {$\bm{v}_1$};
  \node[input] (s-v2) at (10,117) {$\bm{v}_2$};
  \node[compute,minimum width=20mm] (s-f1) at (42,132) {$f_{\bm{\theta}}$};
  \node[compute,minimum width=20mm] (s-f2) at (42,117) {$f_{\bm{\theta}}$};
  \node[compute,minimum width=20mm] (s-g1) at (76,132) {$g_{\bm{\theta}}$};
  \node[compute,minimum width=20mm] (s-g2) at (76,117) {$g_{\bm{\theta}}$};
  \node[output,minimum width=14mm] (s-z1) at (108,132) {$\bm{z}_1$};
  \node[output,minimum width=14mm] (s-z2) at (108,117) {$\bm{z}_2$};
  \node[loss,minimum width=26mm] (s-loss) at (148,124.5) {相似度$\bm{S}$\\与配对损失};
  \draw[forward] (s-v1.east)--(s-f1.west);
  \draw[forward] (s-f1.east)--(s-g1.west);
  \draw[forward] (s-g1.east)--(s-z1.west);
  \draw[forward] (s-v2.east)--(s-f2.west);
  \draw[forward] (s-f2.east)--(s-g2.west);
  \draw[forward] (s-g2.east)--(s-z2.west);
  \draw[NNWire,line width=1pt,rounded corners=1mm] (s-z1.east) -- (126,132) -- (126,124.5);
  \draw[NNWire,line width=1pt,rounded corners=1mm] (s-z2.east) -- (126,117) -- (126,124.5);
  \draw[forward] (126,124.5)--(s-loss.west);
  \draw[gradient] (s-loss.north) -- (148,140) -- (42,140) -- (s-f1.north);
  \draw[gradient] (76,140)--(s-g1.north);
  \draw[gradient] (108,140)--(s-z1.north);
  \draw[gradient] (s-loss.south) -- (148,109) -- (42,109) -- (s-f2.south);
  \draw[gradient] (76,109)--(s-g2.south);
  \draw[gradient] (108,109)--(s-z2.south);

  \node[font=\figurefont\bfseries\footnotesize,text=BookInk,anchor=west] at (5,93) {(b) 跨视图预测：BYOL};
  \node[note,anchor=east] at (163,93) {只画$A\to B$；交换视图另算一项};
  \node[input] (b-va) at (10,78) {$\bm{v}_A$};
  \node[compute,minimum width=20mm] (b-fo) at (42,78) {$f_{\bm{\theta}}$};
  \node[compute,minimum width=20mm] (b-go) at (76,78) {$g_{\bm{\theta}}$};
  \node[compute,minimum width=20mm] (b-qo) at (108,78) {$q_{\bm{\theta}}$};
  \node[input] (b-vb) at (10,59) {$\bm{v}_B$};
  \node[compute,minimum width=20mm] (b-ft) at (42,59) {$f_{\bm{\xi}}$};
  \node[compute,minimum width=20mm] (b-gt) at (76,59) {$g_{\bm{\xi}}$};
  \node[output,minimum width=20mm] (b-sg) at (108,59) {$\operatorname{sg}$};
  \node[loss,minimum width=20mm] (b-loss) at (148,68.5) {$\ell_{A\to B}$};
  \draw[forward] (b-va.east)--(b-fo.west);
  \draw[forward] (b-fo.east)--(b-go.west);
  \draw[forward] (b-go.east)--(b-qo.west);
  \draw[forward,rounded corners=1mm] (b-qo.east) -| (b-loss.north);
  \draw[forward] (b-vb.east)--(b-ft.west);
  \draw[forward] (b-ft.east)--(b-gt.west);
  \draw[forward] (b-gt.east)--(b-sg.west);
  \draw[forward,rounded corners=1mm] (b-sg.east) -| (b-loss.south);
  \draw[gradient] (b-loss.east) -- (164,68.5) -- (164,86) -- (42,86) -- (b-fo.north);
  \draw[gradient] (76,86)--(b-go.north);
  \draw[gradient] (108,86)--(b-qo.north);
  \draw[gradient stop] (164,68.5) -- (164,52) -- (108,52) -- (b-sg.south);
  \draw[ema] (b-fo.south)--(b-ft.north);
  \draw[ema] (b-go.south)--(b-gt.north);
  \node[note,text=BookTeal] at (59,68.5) {EMA};

  \node[font=\figurefont\bfseries\footnotesize,text=BookInk,anchor=west] at (5,37) {(c) 对称匹配：VICReg};
  \node[note,anchor=east] at (163,37) {两路共享编码器与扩展头；无停止梯度};
  \node[input] (v-va) at (10,22) {$\bm{v}_A$};
  \node[compute,minimum width=28mm] (v-fa) at (60,22) {$h_{\bm{\phi}}\circ f_{\bm{\theta}}$};
  \node[output,minimum width=14mm] (v-za) at (100,22) {$\bm{Z}$};
  \node[input] (v-vb) at (10,4) {$\bm{v}_B$};
  \node[compute,minimum width=28mm] (v-fb) at (60,4) {$h_{\bm{\phi}}\circ f_{\bm{\theta}}$};
  \node[output,minimum width=14mm] (v-zb) at (100,4) {$\bm{Z}'$};
  \node[loss,minimum width=34mm] (v-loss) at (148,13) {匹配$s(\bm{Z},\bm{Z}')$\\方差$v(\bm{Z}),v(\bm{Z}')$\\协方差$c(\bm{Z}),c(\bm{Z}')$};
  \draw[forward] (v-va.east)--(v-fa.west);
  \draw[forward] (v-fa.east)--(v-za.west);
  \draw[forward] (v-vb.east)--(v-fb.west);
  \draw[forward] (v-fb.east)--(v-zb.west);
  \draw[NNWire,line width=1pt,rounded corners=1mm] (v-za.east) -- (123,22) -- (123,13);
  \draw[NNWire,line width=1pt,rounded corners=1mm] (v-zb.east) -- (123,4) -- (123,13);
  \draw[forward] (123,13)--(v-loss.west);
  \draw[gradient] (v-loss.north) -- (148,29) -- (60,29) -- (v-fa.north);
  \draw[gradient] (100,29)--(v-za.north);
  \draw[gradient] (v-loss.south) -- (148,-3) -- (60,-3) -- (v-fb.south);
  \draw[gradient] (100,-3)--(v-zb.south);

  \draw[forward] (3,-17) -- (16,-17);
  \node[note,anchor=west] at (18,-17) {前向值};
  \draw[gradient] (69,-17) -- (56,-17);
  \node[note,anchor=west] at (72,-17) {当前步梯度};
  \draw[ema] (120,-17) -- (133,-17);
  \node[note,anchor=west] at (136,-17) {跨步参数平均};
\end{tikzpicture}
